\chapter{Virgule flottante et algèbre linéaire numérique}\label{ch:qm:floating-point-and-numerical-linear-algebra}

En janvier 1982, la Vancouver Stock Exchange lance un indice à 1000. Il est recalculé après chaque transaction, quelques milliers de fois par jour,
et après chaque recalcul, il est tronqué à trois décimales, et non arrondi. Vingt-deux mois plus tard, le vendredi 25 novembre 1983, il clôture
à 524.811 alors que le marché avait en fait progressé~; recalculé sur ce week-end, il s'établit à 1098.892. Chaque troncature coûte moins
d'un millième de point. Un million d'entre elles coûtent la moitié de l'indice. Ce chapitre traite de ce que fait la machine des nombres~: le format
d'un \omterm{def:qm:floating-point-and-numerical-linear-algebra:fp}{nombre à virgule flottante} et l'arrondi de chaque opération, l'annulation et l'accumulation des erreurs, le \omterm{def:qm:floating-point-and-numerical-linear-algebra:cond}{conditionnement} d'un problème et la
stabilité d'un algorithme, les factorisations sur lesquelles repose l'algèbre linéaire numérique, les méthodes itératives de résolution et la reproductibilité
qu'un système de gestion des risques doit à ses utilisateurs.

\section{Représentation}

\begin{definition}[Nombre à virgule flottante, epsilon machine, unité de dernière position]\label{def:qm:floating-point-and-numerical-linear-algebra:fp}
Un \emph{nombre à virgule flottante}\index{nombre à virgule flottante} (floating-point number) binaire est $\pm m \times 2^e$ avec une mantisse entière $m$ de $p$ bits et un exposant $e$ dans un intervalle fixe~; la norme IEEE 754
en double précision (binary64), $p = 53$ et les exposants vont de $-1022$ à $1023$. L'\emph{epsilon machine}\index{epsilon machine} (machine epsilon) est l'écart entre 1 et le nombre immédiatement supérieur, $2^{1-p} =
2^{-52} \approx 2.2 \times 10^{-16}$~; l'unité d'arrondi $u$ en est la moitié (certains textes, et certaines bibliothèques, appellent $u$ lui-même epsilon machine). Une \emph{unité de dernière position}\index{unité de dernière position} (unit in the last place, ulp) de $x$ est l'écart entre les deux
nombres à virgule flottante qui encadrent $x$.
\end{definition}

\begin{definition}[Nombre sous-normal, multiplication-addition fusionnée]\label{def:qm:floating-point-and-numerical-linear-algebra:fma}
Un \emph{nombre sous-normal}\index{nombre sous-normal} (subnormal number) comble l'écart entre zéro et le plus petit nombre normal, par une mantisse réduite, de sorte que $x - y = 0$ seulement si $x = y$. Une
\emph{multiplication-addition fusionnée}\index{multiplication-addition fusionnée} (fused multiply-add, FMA) calcule $a \times b + c$ avec un seul arrondi.
\end{definition}

Le modèle standard de l'arithmétique~: chaque opération élémentaire renvoie le résultat exact arrondi, $\mathrm{fl}(x \circ y) = (x \circ y)(1 + \delta)$ avec $|\delta| \le u = 2^{-53}$. La plupart des fractions décimales
ne sont pas représentables~: $0.1$ est stocké comme le nombre binary64 le plus proche, légèrement supérieur à $0.1$, de sorte que $0.1 \times 10 - 1$ s'évalue exactement à 0 avec deux arrondis, alors qu'une \omterm{def:qm:floating-point-and-numerical-linear-algebra:fma}{multiplication-addition
fusionnée}, qui n'arrondit qu'une fois, renvoie $2^{-54}$. Les deux résultats sont corrects au sens de la norme IEEE. Un compilateur qui contracte $a \times b + c$ en une \omterm{def:qm:floating-point-and-numerical-linear-algebra:fma}{multiplication-addition fusionnée} modifie le résultat~: GCC le fait par
défaut hors d'un mode C conforme aux normes (\texttt{-ffp-contract=fast}) lorsque la cible dispose de l'instruction, ce qui est l'une des raisons pour lesquelles un moteur de risque en C++ et un autre en Python peuvent diverger
sur les derniers bits.

\section{Annulation et accumulation des erreurs}

\begin{definition}[Annulation catastrophique, sommation compensée, algorithme de Welford]\label{def:qm:floating-point-and-numerical-linear-algebra:cancel}
L'\emph{annulation catastrophique}\index{annulation catastrophique} (catastrophic cancellation) est la perte de précision relative lorsque deux nombres arrondis presque égaux sont soustraits~: les premiers chiffres s'annulent et les
erreurs d'arrondi subsistent. La \emph{sommation compensée}\index{sommation compensée} (compensated summation) (Kahan, 1965~; variante de Neumaier, 1974) porte l'erreur d'arrondi de chaque addition dans une variable séparée
et la rajoute. L'\emph{algorithme de Welford}\index{algorithme de Welford} (1962) met à jour la moyenne courante et la somme des carrés des écarts, une observation à la fois.
\end{definition}

\begin{proposition}[Bornes d'erreur de sommation]\label{prop:qm:floating-point-and-numerical-linear-algebra:sum}
La sommation récursive de $n$ nombres a une erreur d'au plus $(n - 1)u\sum_i|x_i| + O(u^2)$~; la sommation par paires d'au plus $\lceil\log_2n\rceil u\sum_i|x_i| + O(u^2)$~; la \omterm{def:qm:floating-point-and-numerical-linear-algebra:cancel}{sommation compensée} d'au plus
$2u\sum_i|x_i| + O(nu^2)\sum_i|x_i|$, indépendante de $n$ au premier ordre.
\end{proposition}

\begin{proof}
Dans la sommation récursive, la somme partielle $s_k$ est arrondie une fois à chaque étape, de sorte que $x_1$ subit $n - 1$ arrondis, chacun apportant une erreur relative d'au plus $u$ sur une somme partielle bornée par
$\sum|x_i|$~; la sommation par paires ne fait passer chaque terme que par $\lceil\log_2n\rceil$ additions. Dans la \omterm{def:qm:floating-point-and-numerical-linear-algebra:cancel}{sommation compensée}, l'erreur de chaque addition est calculée exactement (pour $|s| \ge |x|$,
$(s - t) + x$ est exact) et réinjectée, de sorte que seuls subsistent la correction finale et les termes du second ordre.
\end{proof}

Sur un million d'entrées de P\&L couvrant sept ordres de grandeur, l'erreur par rapport à la somme rationnelle exacte est de $2.5 \times 10^{-6}$ pour la boucle naïve, $2.8 \times 10^{-8}$ pour la sommation par paires
et $2.2 \times 10^{-9}$, soit moins d'une \omterm{def:qm:floating-point-and-numerical-linear-algebra:fp}{unité de dernière position} du total ($7.5 \times 10^{-9}$), pour la sommation de Neumaier (la \cref{fig:qm:floating-point-and-numerical-linear-algebra:sum}). Aucune n'a d'importance pour un rapport de P\&L~; toutes
comptent lorsque deux systèmes doivent concorder au centime près et que l'un d'eux additionne dans un ordre différent.

\begin{omfigure}
\begin{tikzpicture}
\begin{loglogaxis}[width=10cm,height=5cm,xlabel={nombre d'entrées de P\&L},ylabel={erreur absolue},
  tick label style={font=\small},label style={font=\small},xmin=700,xmax=1.5e6,ymin=1e-13,ymax=1e-4,grid=major,grid style={gray!25},
  legend columns=3,legend style={font=\footnotesize,at={(0.5,-0.3)},anchor=north,draw=none,fill=none,
  /tikz/every even column/.append style={column sep=8pt}},table/col sep=comma]
\addplot[omThm,very thick,mark=*,mark size=1.6pt] table[x=n,y=naive]{figdata/methods/25-floating-point-and-numerical-linear-algebra/summation.csv};
\addplot[black,very thick,mark=square*,mark size=1.6pt] table[x=n,y=pairwise]{figdata/methods/25-floating-point-and-numerical-linear-algebra/summation.csv};
\addplot[omDef,very thick,mark=triangle*,mark size=2pt] table[x=n,y=neumaier]{figdata/methods/25-floating-point-and-numerical-linear-algebra/summation.csv};
\legend{boucle naïve,sommation par paires,Neumaier}
\end{loglogaxis}
\end{tikzpicture}

\omcaption{Erreur de trois algorithmes de sommation par rapport à la somme exacte (rationnelle) de $n$ entrées de P\&L simulées dont les ordres de grandeur parcourent cycliquement sept ordres. La somme compensée de Neumaier reste en dessous de
l'\omterm{def:qm:floating-point-and-numerical-linear-algebra:fp}{unité de dernière position} du total. Données~: le tutoriel du chapitre, avec graine fixée.}\label{fig:qm:floating-point-and-numerical-linear-algebra:sum}
\end{omfigure}

La formule classique de la variance, $(\sum x_i^2 - (\sum x_i)^2/n)/(n - 1)$, soustrait deux grands nombres presque égaux. Sur dix mille prix égaux à un décalage plus un tirage uniforme, son erreur relative
vaut $3 \times 10^{-13}$ avec un décalage de 1, $8.0 \times 10^{-7}$ avec un décalage de $10^4$ (prix proches de 10\,000), 1.3~\% avec un décalage de $10^6$, et avec un décalage de $10^8$ elle renvoie 86.8 pour une variance vraie de 0.082.
La formule en deux passes, qui soustrait d'abord la moyenne, reste voisine de $10^{-12}$ ou en deçà~; la mise à jour de Welford reste en deçà de $1.6 \times 10^{-8}$ même à $10^8$
(la \cref{fig:qm:floating-point-and-numerical-linear-algebra:var}).

\begin{omfigure}
\begin{tikzpicture}
\begin{loglogaxis}[width=10cm,height=5cm,xlabel={décalage des données},ylabel={erreur relative de la variance},
  tick label style={font=\small},label style={font=\small},xmin=0.5,xmax=2e8,ymin=1e-17,ymax=1e4,grid=major,grid style={gray!25},
  legend columns=3,legend style={font=\footnotesize,at={(0.5,-0.3)},anchor=north,draw=none,fill=none,
  /tikz/every even column/.append style={column sep=8pt}},table/col sep=comma]
\addplot[omThm,very thick,mark=*,mark size=1.6pt] table[x=offset,y=textbook]{figdata/methods/25-floating-point-and-numerical-linear-algebra/variance.csv};
\addplot[omDef,very thick,mark=triangle*,mark size=2pt] table[x=offset,y=welford]{figdata/methods/25-floating-point-and-numerical-linear-algebra/variance.csv};
\addplot[black,very thick,mark=square*,mark size=1.6pt] table[x=offset,y=twopass]{figdata/methods/25-floating-point-and-numerical-linear-algebra/variance.csv};
\legend{formule classique en une passe,Welford,deux passes}
\end{loglogaxis}
\end{tikzpicture}

\omcaption{Erreur relative de trois formules de variance sur 10\,000 valeurs égales à un décalage plus un tirage uniforme sur $(0, 1)$ (variance vraie 0.082), par rapport au résultat rationnel exact. La
formule classique en une passe perd toute précision à mesure que le décalage croît. Données~: le tutoriel du chapitre, avec graine fixée.}\label{fig:qm:floating-point-and-numerical-linear-algebra:var}
\end{omfigure}

L'indice de Vancouver obéit au même mécanisme, avec un biais. L'arrondi à trois décimales produit des erreurs de l'un ou l'autre signe qui se compensent~; la troncature retire, en moyenne, un demi-millième à chaque
recalcul, toujours à la baisse. À 2\,400 recalculs par jour sur 480 jours de cotation, la perte attendue vaut $0.0005 \times 1\,152\,000 = 576$ points, contre un écart documenté de $1098.892 -
524.811 = 574$~; au chiffre « environ 3\,000 » par jour parfois cité, on obtiendrait 720. Une simulation des recalculs, avec un indice exact qui monte jusqu'à 1098.9, aboutit à 422 lorsqu'il est tronqué
(la dérive se cumule aux mouvements propres de l'indice) et à 1098.57 lorsqu'il est arrondi (la \cref{fig:qm:floating-point-and-numerical-linear-algebra:vancouver}).

\begin{omfigure}
\begin{tikzpicture}
\begin{axis}[width=11cm,height=5cm,xlabel={jour de cotation depuis janvier 1982},ylabel={indice},
  tick label style={font=\small},label style={font=\small},xmin=0,xmax=480,ymin=300,ymax=1450,grid=major,grid style={gray!25},legend columns=2,legend cell align=left,
  legend style={font=\footnotesize,at={(0.03,0.97)},anchor=north west,fill=white,draw=none,/tikz/every even column/.append style={column sep=8pt}},table/col sep=comma]
\addplot[omDef,very thick] table[x=day,y=exact]{figdata/methods/25-floating-point-and-numerical-linear-algebra/vancouver.csv};
\addplot[omThm,very thick] table[x=day,y=truncated]{figdata/methods/25-floating-point-and-numerical-linear-algebra/vancouver.csv};
\legend{exact,tronqué à trois décimales}
\end{axis}
\end{tikzpicture}

\omcaption{Un indice simulé recalculé 2\,400 fois par jour pendant 480 jours de cotation, exactement et avec troncature à trois décimales après chaque recalcul~; l'indice exact est calibré pour finir à la
valeur corrigée de Vancouver, 1098.892. Données~: le tutoriel du chapitre, avec graine fixée.}\label{fig:qm:floating-point-and-numerical-linear-algebra:vancouver}
\end{omfigure}

\section{Conditionnement et stabilité}

\begin{definition}[Conditionnement, erreur inverse, stabilité inverse]\label{def:qm:floating-point-and-numerical-linear-algebra:cond}
Le \emph{conditionnement}\index{conditionnement} d'un problème majore la variation relative de sa réponse en fonction de la variation relative de ses données~; pour la résolution de $Ax = b$, il vaut $\kappa(A) =
\lVert A\rVert\lVert A^{-1}\rVert$, le rapport des valeurs singulières extrêmes pour la norme 2. L'\emph{erreur inverse}\index{erreur inverse} d'une réponse calculée est la plus petite perturbation des données
pour laquelle elle est exacte. Un algorithme a la \emph{stabilité inverse}\index{stabilité inverse} si son erreur inverse est toujours de l'ordre de l'unité d'arrondi.
\end{definition}

\begin{proposition}[Erreur directe, conditionnement et erreur inverse]\label{prop:qm:floating-point-and-numerical-linear-algebra:fwd}
Si $\hat x$ résout $(A + \Delta A)\hat x = b$ avec $\lVert\Delta A\rVert/\lVert A\rVert = \epsilon$ et $\kappa(A)\epsilon < 1$, alors $\lVert\hat x - x\rVert/\lVert x\rVert \le \kappa(A)\epsilon/(1 - \kappa(A)\epsilon)$~:
au premier ordre, l'erreur directe est majorée par le produit du \omterm{def:qm:floating-point-and-numerical-linear-algebra:cond}{conditionnement} et de l'\omterm{def:qm:floating-point-and-numerical-linear-algebra:cond}{erreur inverse}.
\end{proposition}

\begin{proof}
Retrancher $Ax = b$ donne $A(\hat x - x) = -\Delta A\,\hat x$, donc $\lVert\hat x - x\rVert \le \lVert A^{-1}\rVert\lVert\Delta A\rVert\lVert\hat x\rVert = \kappa\epsilon\lVert\hat x\rVert$, et $\lVert\hat x\rVert \le
\lVert x\rVert + \lVert\hat x - x\rVert$ donne la majoration.
\end{proof}

Un algorithme à \omterm{def:qm:floating-point-and-numerical-linear-algebra:cond}{stabilité inverse} sur un problème de \omterm{def:qm:floating-point-and-numerical-linear-algebra:cond}{conditionnement} $10^8$ peut perdre huit de ses seize chiffres,
et aucun algorithme ne peut faire mieux, car les données ne sont connues qu'à l'arrondi près. Ce qu'un algorithme peut faire, c'est ne pas dégrader le \omterm{def:qm:floating-point-and-numerical-linear-algebra:cond}{conditionnement}. Les équations normales des moindres carrés,
$X^\top X\beta = X^\top y$, résolvent un problème de \omterm{def:qm:floating-point-and-numerical-linear-algebra:cond}{conditionnement} $\kappa(X)^2$~; une \omterm{def:qm:floating-point-and-numerical-linear-algebra:factor}{factorisation QR} de $X$ en résout un de \omterm{def:qm:floating-point-and-numerical-linear-algebra:cond}{conditionnement} $\kappa(X)$. Avec cinq régresseurs de \omterm{def:qm:floating-point-and-numerical-linear-algebra:cond}{conditionnement} $10^7$, les
équations normales perdent tous les chiffres sauf deux des coefficients (erreur relative $5 \times 10^{-3}$), la solution QR en conserve une dizaine (erreur voisine de $10^{-10}$, dont les chiffres exacts relèvent de l'arrondi et varient
avec la bibliothèque d'algèbre linéaire et le processeur)
(\cref{fig:qm:floating-point-and-numerical-linear-algebra:normal}). C'est la \omterm{def:qm:linear-models-under-stress:vif}{multicolinéarité} du chapitre~16 vue du côté de la machine.

\begin{omfigure}
\begin{tikzpicture}
\begin{loglogaxis}[width=10cm,height=5cm,xlabel={conditionnement de la matrice de plan $\kappa(X)$},ylabel={erreur relative de $\hat\beta$},
  tick label style={font=\small},label style={font=\small},xmin=5,xmax=2e7,ymin=1e-16,ymax=1e-1,grid=major,grid style={gray!25},
  legend cell align=left,legend style={font=\footnotesize,at={(0.03,0.97)},anchor=north west,fill=white,draw=none},table/col sep=comma]
\addplot[omThm,very thick,mark=*,mark size=1.6pt] table[x=cond,y=normal]{figdata/methods/25-floating-point-and-numerical-linear-algebra/normal.csv};
\addplot[omDef,very thick,mark=square*,mark size=1.6pt] table[x=cond,y=qr]{figdata/methods/25-floating-point-and-numerical-linear-algebra/normal.csv};
\legend{équations normales,QR}
\end{loglogaxis}
\end{tikzpicture}

\omcaption{Moindres carrés sur un système exactement compatible de 500 observations et cinq régresseurs dont la matrice de plan a pour \omterm{def:qm:floating-point-and-numerical-linear-algebra:cond}{conditionnement} $\kappa$~: erreur relative des coefficients obtenue par les
équations normales (qui élèvent $\kappa$ au carré) et par une \omterm{def:qm:floating-point-and-numerical-linear-algebra:factor}{factorisation QR}. Données~: le tutoriel du chapitre, avec graine fixée.}\label{fig:qm:floating-point-and-numerical-linear-algebra:normal}
\end{omfigure}

\section{Factorisations}

\begin{definition}[LU, Cholesky, QR et décomposition en valeurs singulières]\label{def:qm:floating-point-and-numerical-linear-algebra:factor}
La \emph{factorisation LU}\index{factorisation LU} écrit $PA = LU$ avec une permutation $P$ (pivotage partiel), $L$ triangulaire inférieure à diagonale unité et $U$ triangulaire supérieure. La
\emph{factorisation de Cholesky}\index{factorisation de Cholesky} d'une matrice symétrique définie positive est $A = LL^\top$ avec $L$ triangulaire inférieure. La \emph{factorisation QR}\index{factorisation QR} écrit $A = QR$ avec
des colonnes orthonormées dans $Q$ et $R$ triangulaire supérieure. La \emph{décomposition en valeurs singulières}\index{décomposition en valeurs singulières} écrit $A = U\Sigma V^\top$ avec $U$ et $V$ orthonormées et
$\Sigma$ diagonale à coefficients positifs ou nuls. L'\emph{algorithme de Thomas (matrices tridiagonales)}\index{algorithme de Thomas (matrices tridiagonales)} résout un système tridiagonal par une élimination avant et une
substitution arrière, en $O(n)$ opérations.
\end{definition}

\begin{proposition}[La factorisation de Cholesky réussit exactement pour les matrices définies positives]\label{prop:qm:floating-point-and-numerical-linear-algebra:chol}
La récurrence de Cholesky $\ell_{jj} = (a_{jj} - \sum_{k<j}\ell_{jk}^2)^{1/2}$, $\ell_{ij} = (a_{ij} - \sum_{k<j}\ell_{ik}\ell_{jk})/\ell_{jj}$ trouve un pivot strictement positif à chaque étape si et seulement si $A$ est
symétrique définie positive.
\end{proposition}

\begin{proof}
Le $j$-ième pivot est le quotient des mineurs principaux dominants d'ordres $j$ et $j - 1$, donc tous les pivots sont positifs si et seulement si tous les mineurs dominants le sont, ce qui, par le critère de
Sylvester, équivaut à la définie positivité.
\end{proof}

Un échec de la \omterm{def:qm:floating-point-and-numerical-linear-algebra:factor}{factorisation de Cholesky} est donc un test, et un test utile~: la matrice de corrélation par paires du chapitre~23 (cinquante séries, la moitié des données manquantes) échoue au sixième pivot, qui vaut $-0.14$. Ajouter un
multiple de l'identité («~jitter~») ne peut réussir tant que le multiple ne dépasse pas, en valeur absolue, la valeur propre la plus négative, 1.42, soit plus que la diagonale unité elle-même, ce qui détruit les corrélations~; la \omterm{def:qm:convex-optimisation:ncm}{matrice de corrélation la plus proche} du chapitre~23 passe avec un jitter de
$10^{-8}$. Chaque factorisation a son usage~: la \omterm{def:qm:floating-point-and-numerical-linear-algebra:factor}{factorisation LU} pour les systèmes généraux, celle de Cholesky pour les covariances (et pour simuler des variables gaussiennes corrélées), la \omterm{def:qm:floating-point-and-numerical-linear-algebra:factor}{factorisation QR} pour les moindres carrés, la \omterm{def:qm:floating-point-and-numerical-linear-algebra:factor}{décomposition en valeurs singulières} pour le rang et le \omterm{def:qm:floating-point-and-numerical-linear-algebra:cond}{conditionnement}, et
l'algorithme de Thomas pour les grilles aux différences finies du chapitre~27.

\section{Solveurs itératifs}

\begin{definition}[Méthode du gradient conjugué, préconditionneur]\label{def:qm:floating-point-and-numerical-linear-algebra:cg}
La \emph{méthode du gradient conjugué}\index{méthode du gradient conjugué} (Hestenes et Stiefel, 1952) résout $Ax = b$ pour $A$ symétrique définie positive en minimisant $\frac12x^\top Ax - b^\top x$ le long
de directions mutuellement $A$-conjuguées construites à partir des résidus. Un \emph{préconditionneur}\index{préconditionneur} $M \approx A$ remplace le système par un système en $M^{-1}A$, de \omterm{def:qm:floating-point-and-numerical-linear-algebra:cond}{conditionnement} plus petit.
\end{definition}

En arithmétique exacte, la \omterm{def:qm:floating-point-and-numerical-linear-algebra:cg}{méthode du gradient conjugué} se termine en au plus $n$ étapes, et l'erreur après $k$ étapes décroît au moins comme $2\bigl((\sqrt\kappa - 1)/(\sqrt\kappa + 1)\bigr)^k$. En virgule flottante, les
directions perdent leur conjugaison et la terminaison finie est perdue, mais le taux de convergence subsiste. Sur un système de dimension 400 et de \omterm{def:qm:floating-point-and-numerical-linear-algebra:cond}{conditionnement} $1.1 \times 10^5$, dont les lignes et les colonnes sont mal
mises à l'échelle, la \omterm{def:qm:floating-point-and-numerical-linear-algebra:cg}{méthode du gradient conjugué} a besoin d'environ 1\,900 itérations pour réduire le résidu d'un facteur $10^8$, soit près de cinq fois $n$ (le nombre exact varie de quelques dizaines selon l'ordre dans lequel la bibliothèque d'algèbre linéaire effectue ses sommes, un seul thread ou quatre~: la perte de conjugaison est un arrondi, et l'arrondi dépend de cet ordre)~; avec un \omterm{def:qm:floating-point-and-numerical-linear-algebra:cg}{préconditionneur} de Jacobi (diagonal), qui ramène le \omterm{def:qm:floating-point-and-numerical-linear-algebra:cond}{conditionnement} à environ
100, elle en demande 102 (\cref{fig:qm:floating-point-and-numerical-linear-algebra:cg}).

\begin{omfigure}
\begin{tikzpicture}
\begin{semilogyaxis}[width=10cm,height=5cm,xlabel={itération},ylabel={résidu relatif},
  tick label style={font=\small},label style={font=\small},xmin=0,xmax=1950,ymin=1e-9,ymax=100,grid=major,grid style={gray!25},xtick={0,500,1000,1500},unbounded coords=jump,
  /pgf/number format/1000 sep={\,},legend cell align=left,legend style={font=\footnotesize,at={(0.97,0.97)},anchor=north east,fill=white,draw=none},table/col sep=comma]
\addplot[omThm,thick] table[x=iter,y=plain]{figdata/methods/25-floating-point-and-numerical-linear-algebra/cg.csv};
\addplot[omDef,very thick] table[x=iter,y=pre]{figdata/methods/25-floating-point-and-numerical-linear-algebra/cg.csv};
\legend{gradient conjugué,préconditionné de Jacobi}
\end{semilogyaxis}
\end{tikzpicture}

\omcaption{Gradient conjugué sur un système symétrique défini positif de dimension 400, mal mis à l'échelle (\omterm{def:qm:floating-point-and-numerical-linear-algebra:cond}{conditionnement} $1.1 \times 10^5$), avec et sans \omterm{def:qm:floating-point-and-numerical-linear-algebra:cg}{préconditionneur} diagonal~: résidu relatif
en fonction de l'itération (la série préconditionnée s'arrête à 102). Données~: le tutoriel du chapitre, avec graine fixée.}\label{fig:qm:floating-point-and-numerical-linear-algebra:cg}
\end{omfigure}

\section{Les bogues qui coûtent cher}

\begin{definition}[Reproductibilité bit à bit]\label{def:qm:floating-point-and-numerical-linear-algebra:repro}
Un calcul a la \emph{reproductibilité bit à bit}\index{reproductibilité bit à bit} si les mêmes entrées donnent les mêmes bits à chaque exécution, sur chaque machine, avec chaque compilateur et chaque langage employé pour les calculer.
\end{definition}

L'addition en virgule flottante n'est pas associative, si bien que tout ce qui modifie l'ordre des opérations modifie les derniers bits~: une réduction parallèle avec un nombre de threads différent, une boucle vectorisée, un
compilateur qui contracte en multiplications-additions fusionnées (FMA), une bibliothèque qui somme par blocs. Ces bits se propagent alors à travers des seuils (une limite franchie ou non, une transaction exécutée ou non) et deviennent de véritables
écarts. Le kit interne de la société fixe donc l'ordre de chaque opération~: ses trois jumeaux, en Python, C++20 et Rust, effectuent les mêmes opérations dans le même ordre sur les mêmes entrées SplitMix64
et s'accordent bit à bit sur les sommations naïve, par paires et compensée de dix mille écritures de P\&L, sur la moyenne et la variance calculées par l'\omterm{def:qm:floating-point-and-numerical-linear-algebra:cancel}{algorithme de Welford} pour dix mille cours, et sur un log-sum-exp et une résolution tridiagonale~; un
test dans chaque langage vérifie les mêmes motifs de 64~bits. Le même test avec une \omterm{def:qm:floating-point-and-numerical-linear-algebra:fma}{multiplication-addition fusionnée} à la place d'une multiplication et d'une addition montre la différence~: $0.1 \times 10 - 1$ vaut 0 dans un cas et $2^{-54}$
dans l'autre. (Avec \texttt{g++ -std=c++20 -O2} sur x86-64, aucune contraction n'a lieu, car le jeu d'instructions de base ne comporte pas de \omterm{def:qm:floating-point-and-numerical-linear-algebra:fma}{multiplication-addition fusionnée}~; \texttt{-march=native} sur un processeur
récent l'active.)

Le mécanisme de Vancouver, un arrondi systématique répété des millions de fois, a coûté plus que de l'argent. Le 25 février 1991, une batterie Patriot de défense antiaérienne à Dhahran, en Arabie saoudite, n'a pas réussi à
intercepter un missile Scud, qui a frappé une caserne de l'armée américaine et tué 28 Américains. Le General Accounting Office a imputé la défaillance à l'horloge du système~: le temps était compté en dixièmes de seconde et
multiplié par $\frac1{10}$ dans des registres de 24~bits, et après environ 100 heures de fonctionnement de la batterie, le temps calculé était décalé de 0.3433 seconde, suffisant pour déplacer la fenêtre de poursuite de 687~mètres. Les
chiffres du rapport correspondent à ce que donne $\frac1{10}$ tronqué à 23 positions binaires~: il est inférieur de $9.5 \times 10^{-8}$, et $3\,600\,000$ dixièmes de seconde accumulent une dérive de $0.3433$ seconde. Des données israéliennes avaient montré
la perte de précision au bout de huit heures~; le logiciel corrigé est arrivé le lendemain de l'attaque.

Tous les bogues numériques ne se logent pas dans le dernier bit. Le groupe de travail de direction qui a examiné les pertes de 2012 de la Chief Investment Office de JPMorgan Chase a signalé, parmi les erreurs relevées dans un nouveau modèle de valeur en risque
exécuté dans des tableurs, que «~après avoir soustrait l'ancien taux du nouveau, le tableur a divisé par leur somme au lieu de leur moyenne~», ce qui «~a probablement eu pour effet
d'atténuer la volatilité d'un facteur deux et de faire baisser la VaR~». La parade est la même que pour l'arrondi~: une implémentation indépendante, des résultats de référence vérifiés par des tests, et un relecteur qui
recalcule à la main un nombre.

\section{Tutoriel~: l'indice qui a perdu la moitié de sa valeur}\label{tut:qm:floating-point-and-numerical-linear-algebra:vancouver}

\begin{tutorial}
\textbf{Objectif.} Reproduire la dérive de Vancouver, mesurer les erreurs de sommation et de variance par rapport à l'arithmétique exacte, observer l'échec de Cholesky et vérifier que trois langages s'accordent bit à bit.
\textbf{État final~:} les \cref{fig:qm:floating-point-and-numerical-linear-algebra:sum,fig:qm:floating-point-and-numerical-linear-algebra:var,fig:qm:floating-point-and-numerical-linear-algebra:vancouver}
et les motifs de bits de \texttt{reference\_results()} dans les trois suites de tests.
\begin{tutsteps}
\item \textbf{\omterm{def:qm:floating-point-and-numerical-linear-algebra:cancel}{Sommation compensée} et mise à jour de Welford} en Python.
\omcode{code/firm/fpkit/firm_fpkit.py}{73}{95}{Sommation de Neumaier et algorithme de Welford (Python).}\label{lst:qm:floating-point-and-numerical-linear-algebra:py}
\item \textbf{Les mêmes opérations en Rust}, dans le même ordre.
\omcode{code/firm/fpkit/rust/src/lib.rs}{54}{78}{Sommation de Neumaier et algorithme de Welford (jumeau Rust).}\label{lst:qm:floating-point-and-numerical-linear-algebra:rs}
\item \textbf{Exécuter} \texttt{vancouver()}, \texttt{summation\_errors()}, \texttt{variance\_errors()}, \texttt{normal\_equations()}, \texttt{cholesky\_on\_pairwise()} et \texttt{cg\_demo()} dans
\texttt{qm\_float.py}~; puis les trois suites de tests de \texttt{firm.fpkit}, ainsi que \texttt{fig\_float.py}.
\end{tutsteps}
\textbf{Pour aller plus loin.} Compiler le jumeau C++ avec \texttt{-march=native -ffp-contract=fast} et observer quels motifs de bits de référence ne tiennent plus~; sommer les P\&L avec la fonction \texttt{sum} de NumPy et expliquer pourquoi
elle est plus proche de la sommation par paires que de la sommation naïve.
\end{tutorial}

\section{Construction~: le kit de calcul numérique}\label{bld:qm:floating-point-and-numerical-linear-algebra:fpkit}

\begin{build}
\bfield{Rôle} Primitives de calcul numérique rigoureuses, identiques dans les trois langages de la mini-société, afin qu'un nombre de risque calculé dans le carnet de recherche, le moteur C++ et
la passerelle Rust soit le même nombre.
\bfield{Interface} \texttt{SplitMix64}~; \texttt{naive\_sum}, \texttt{pairwise\_sum}, \texttt{neumaier\_sum}~; \texttt{welford}, \texttt{welford\_cov}~; \texttt{logsumexp}~; \texttt{thomas}~; en Python également
\texttt{cholesky(A, jitter)}, \texttt{cg(A, b, precond)}, \texttt{fma\_exact}, \texttt{bits}.
\bfield{Règles} Ni fast-math, ni contraction en virgule flottante dans les builds qui doivent reproduire~; ordre de sommation fixe~; chaque jumeau vérifie les motifs de bits de référence~; un Cholesky en échec est signalé avec son
indice, jamais silencieusement corrigé par ajout de jitter.
\bfield{Tests d'acceptation} \texttt{code/firm/fpkit/tests/}, \texttt{cpp/firm\_fpkit\_test.cpp}, \texttt{rust/src/lib.rs}~: bits identiques dans les trois langages~; la sommation de Neumaier égale ici la somme exacte~;
Welford comparé à l'arithmétique rationnelle exacte~; Thomas comparé à une résolution dense~; Cholesky sur des matrices définies positives et indéfinies~; gradient conjugué~; l'exemple de la \omterm{def:qm:floating-point-and-numerical-linear-algebra:fma}{multiplication-addition fusionnée}.
\bfield{Pour aller plus loin} Une somme parallèle reproductible (blocs de taille fixe et compensation)~; arithmétique double-double~; un Cholesky avec pivotage qui renvoie le rang.
\end{build}

\begin{omsources}
\item IEEE Standard for Floating-Point Arithmetic, IEEE Std 754-2019.
\item W.~Kahan, «~Further remarks on reducing truncation errors~», \emph{Communications of the ACM} 8, 1965~; A.~Neumaier, \emph{Zeitschrift für Angewandte Mathematik und Mechanik} 54, 1974.
\item B.~P.~Welford, «~Note on a method for calculating corrected sums of squares and products~», \emph{Technometrics} 4, 1962.
\item M.~R.~Hestenes and E.~Stiefel, «~Methods of conjugate gradients for solving linear systems~», \emph{Journal of Research of the National Bureau of Standards} 49, 1952.
\item N.~J.~Higham, \emph{Accuracy and Stability of Numerical Algorithms}, 2\up{e} éd., SIAM, 2002.
\item D.~Goldberg, «~What every computer scientist should know about floating-point arithmetic~», \emph{ACM Computing Surveys} 23, 1991.
\item Manuel de GCC, «~Options that control optimization~» (\texttt{-ffp-contract}), consulté le 24 septembre 2026.
\item US General Accounting Office, \emph{Patriot Missile Defense: Software Problem Led to System Failure at Dhahran, Saudi Arabia}, GAO/IMTEC-92-26, 4 février 1992.
\item Sur l'indice de Vancouver~: K.~Quinn, \emph{The Wall Street Journal}, 8 novembre 1983, p.~37, et W.~Lilley, \emph{The Toronto Star}, 29 novembre 1983, p.~35, tel que cité dans l'article Wikipedia
«~Vancouver Stock Exchange~» (consulté le 24 septembre 2026).
\end{omsources}

\section{Exercices}

\begin{exercise}[$\star$]\label{exo:qm:floating-point-and-numerical-linear-algebra:1}
Donner l'unité d'arrondi et l'\omterm{def:qm:floating-point-and-numerical-linear-algebra:fp}{epsilon machine} de binary64, ainsi que le nombre de chiffres décimaux significatifs qu'il porte.
\end{exercise}

\begin{exercise}[$\star$]\label{exo:qm:floating-point-and-numerical-linear-algebra:2}
Un indice est tronqué à deux décimales après chacun des 10\,000 recalculs quotidiens. Déterminer la dérive à attendre sur une année de 250 jours.
\end{exercise}

\begin{exercise}[$\star$]\label{exo:qm:floating-point-and-numerical-linear-algebra:3}
Expliquer pourquoi $10^8 + 1 - 10^8$ est exact en double précision, mais pas $10^{16} + 1 - 10^{16}$.
\end{exercise}

\begin{exercise}[$\star\star$]\label{exo:qm:floating-point-and-numerical-linear-algebra:4}
Montrer que l'erreur absolue de la formule classique de la variance est d'ordre $u\sum x_i^2/(n - 1)$, et estimer l'erreur relative pour des cours voisins de 10\,000 avec une variance de 0.08.
\end{exercise}

\begin{exercise}[$\star\star$]\label{exo:qm:floating-point-and-numerical-linear-algebra:5}
Estimer le nombre de chiffres que peuvent perdre les équations normales et la \omterm{def:qm:floating-point-and-numerical-linear-algebra:factor}{factorisation QR} pour une matrice de design de \omterm{def:qm:floating-point-and-numerical-linear-algebra:cond}{conditionnement} $10^5$.
\end{exercise}

\begin{exercise}[$\star\star$]\label{exo:qm:floating-point-and-numerical-linear-algebra:6}
Montrer que l'astuce log-sum-exp, $m + \log\sum e^{x_i - m}$ avec $m = \max x_i$, ne déborde jamais, et calculer $\log(e^{1000} + e^{1000})$.
\end{exercise}

\begin{exercise}[$\star\star\star$]\label{exo:qm:floating-point-and-numerical-linear-algebra:7}
\emph{Programmation.} Sommer un million d'écritures de P\&L simulées de façon naïve, par paires et avec la \omterm{def:qm:floating-point-and-numerical-linear-algebra:cancel}{sommation compensée} de Neumaier, et comparer avec la somme rationnelle exacte.
\end{exercise}

\begin{exercise}[$\star\star\star$]\label{exo:qm:floating-point-and-numerical-linear-algebra:8}
\emph{Repérer le défaut.} «~Nos totaux de risque diffèrent entre le moteur C++ et les rapports Python au dixième chiffre significatif, donc l'un des deux contient un bogue.~»
\end{exercise}

\section{Problème~: l'indice qui a perdu la moitié de sa valeur}

\begin{problem}[{Problème du week-end --- troncature, un million de fois}]\label{pb:qm:floating-point-and-numerical-linear-algebra:1}
Un indice part de 1000 et est recalculé après chaque transaction, 2\,400 fois par jour pendant 480 jours de cotation~; après chaque recalcul, il est tronqué à trois décimales. Le cas documenté~: 524.811 avant la
correction, 1098.892 après.

\medskip\noindent\textbf{Partie I --- La dérive.}
\begin{enumerate}
\item Déterminer la perte espérée due à une troncature, et en expliquer la raison.
\item Calculer la perte totale espérée, et la comparer au déficit documenté.
\item Calculer la perte totale pour 3\,000 recalculs par jour.
\item Indiquer ce que change un arrondi au lieu d'une troncature.
\item Déterminer la valeur finale de l'indice tronqué simulé, et expliquer pourquoi elle est inférieure à celle de la formule.
\end{enumerate}

\medskip\noindent\textbf{Partie II --- Erreur accumulée.}
\begin{enumerate}[resume]
\item Évaluer l'ordre de grandeur des erreurs de sommation naïve, par paires et compensée sur un million d'écritures de P\&L.
\item Évaluer l'erreur de la variance classique pour des cours voisins de 10\,000, $10^6$ et $10^8$.
\item Expliquer pourquoi les formules à deux passes et celle de Welford résistent.
\item Déterminer le nombre de chiffres perdus par les équations normales pour un \omterm{def:qm:floating-point-and-numerical-linear-algebra:cond}{conditionnement} de $10^7$, et par la \omterm{def:qm:floating-point-and-numerical-linear-algebra:factor}{factorisation QR}.
\item Indiquer où et pourquoi la \omterm{def:qm:floating-point-and-numerical-linear-algebra:factor}{factorisation de Cholesky} de la matrice de corrélation par paires échoue.
\end{enumerate}

\medskip\noindent\textbf{Partie III --- Reproductibilité.}
\begin{enumerate}[resume]
\item Indiquer sur quels points les trois jumeaux s'accordent, et comment on le vérifie.
\item Indiquer ce que change une \omterm{def:qm:floating-point-and-numerical-linear-algebra:fma}{multiplication-addition fusionnée} dans $0.1 \times 10 - 1$.
\item Indiquer quelle option de compilation peut introduire silencieusement des multiplications-additions fusionnées.
\item Calculer le nombre d'itérations de la \omterm{def:qm:floating-point-and-numerical-linear-algebra:cg}{méthode du gradient conjugué} nécessaires pour le système mal mis à l'échelle, avec et sans \omterm{def:qm:floating-point-and-numerical-linear-algebra:cg}{préconditionneur} de Jacobi.
\item Expliquer pourquoi la \omterm{def:qm:floating-point-and-numerical-linear-algebra:cg}{méthode du gradient conjugué} ne s'arrête pas après $n$ étapes ici.
\end{enumerate}

\medskip\noindent\textbf{Partie IV --- Jugement.}
\begin{enumerate}[resume]
\item Indiquer comment on aurait détecté l'erreur de Vancouver en une semaine~?
\item Indiquer ce que doit consigner un système de gestion des risques pour que deux exécutions puissent être comparées bit à bit~?
\item Indiquer quand la \omterm{def:qm:floating-point-and-numerical-linear-algebra:repro}{reproductibilité bit à bit} vaut son coût.
\item Énoncer le \emph{résultat nommé}~: la dérive espérée par recalcul multipliée par le nombre de recalculs, à comparer au déficit documenté.
\item En une phrase, indiquer ce que fait la machine à chaque nombre calculé~?
\end{enumerate}
\end{problem}

\section{Questions d'entretien}

\begin{interviewq}[$\star$ \iqroles{développeur, chercheur}]\label{iq:qm:floating-point-and-numerical-linear-algebra:1}
Pourquoi \texttt{0.1 + 0.2 == 0.3} renvoie-t-il faux, et comment comparez-vous des \omterm{def:qm:floating-point-and-numerical-linear-algebra:fp}{nombres à virgule flottante}~?
\end{interviewq}

\begin{interviewq}[$\star\star$ \iqroles{développeur}]\label{iq:qm:floating-point-and-numerical-linear-algebra:2}
Vos moteurs de risque en C++ et en Python divergent sur les derniers chiffres. Énumérez les causes probables.
\end{interviewq}

\begin{interviewq}[$\star\star$ \iqroles{chercheur, développeur}]\label{iq:qm:floating-point-and-numerical-linear-algebra:3}
Comment calculez-vous la variance des prix dans un système en flux continu, et pourquoi ne pas utiliser la somme des carrés~?
\end{interviewq}

\begin{interviewq}[$\star\star$ \iqroles{chercheur}]\label{iq:qm:floating-point-and-numerical-linear-algebra:4}
Pourquoi résoudre les moindres carrés par QR plutôt que par les équations normales~?
\end{interviewq}

\begin{interviewq}[$\star\star$ \iqroles{développeur, chercheur}]\label{iq:qm:floating-point-and-numerical-linear-algebra:5}
Votre \omterm{def:qm:floating-point-and-numerical-linear-algebra:factor}{factorisation de Cholesky} d'une \omterm{def:qm:covariance-estimation-and-random-matrices:cov}{matrice de covariance} échoue. Que faites-vous~?
\end{interviewq}

\begin{interviewq}[$\star\star\star$ \iqroles{développeur}]\label{iq:qm:floating-point-and-numerical-linear-algebra:6}
Comment rendriez-vous une somme parallèle reproductible quel que soit le nombre de threads~?
\end{interviewq}
